\documentclass[10pt]{article}

\usepackage[utf8]{inputenc}

\usepackage[T1]{fontenc}

\usepackage{amsmath}

\usepackage{amsfonts}

\usepackage{amssymb}

\usepackage[version=4]{mhchem}

\usepackage{stmaryrd}

\usepackage{bbold}

\usepackage{graphicx}

\usepackage[export]{adjustbox}

\graphicspath{ {./images/} }



\begin{document}

одной орбите $G L^{k}(n) w_{0}$ групшы $G L^{k}(n)$. Тогда $P=S^{-1}\left(w_{0}\right)$ есть G-С., где $G$ - стабилизатор точки $w_{0}$ в $G L^{k}(n)$, и $S$ является ее тензором Бернара. Напр., риманова метрика задает $O(n)$-структуру, почти симплектич. структура - $S p\left(\frac{n}{2}, \mathbb{R}\right)$-структуру, почти комплексная структура - $G L\left(\frac{n}{2}, \mathbb{C}\right)$-структуру, аффинная связность без кручения - $G L(n)$-структуру 2-го порядка $(G L(n)$ рассматривается здесь как подгруппа группы $\left.G L^{2}(n)\right)$. Афффинор (поле эндоморфизмов) задает $G$-С. тогда и только тогда, когда он во всех точках приводится к одной и той же жордановой нормальной форме $A$, причем $G$ есть централизатор матрицы $A$ в $G L(n)$.



Элементы многообразия $M_{k}$ можно рассматривать как кореперы порядка 1 на $M_{k-1}$, что позволяет рассматривать естественное расслоение $\pi^{k}: M_{k} \rightarrow M_{k-1}$ как $N^{k}$-структуру 1-го порядка, где $N^{k}$ - ядро естественного гомоморфизма $G L^{k}(n) \rightarrow G L^{k-1}(n)$. С каждой $G$-С. $\pi: P \rightarrow M$ порядка $k$ связывается послөдовательность $G$-С. 1-го порядка



\[

P \longrightarrow P_{-1} \longrightarrow P_{-2} \longrightarrow \ldots \longrightarrow P_{-k}=M

\]



где $P_{-i}=\pi^{k}\left(P_{-i+1}\right) \subset M_{k-i} \cdot$ Благодаря әтому изучение $G-C$. высших порядков сводится к изучению $G-C$. первого порядка. Корепер $u_{x}^{1} \in M_{1}$ можно рассматривать как изоморфизм $u_{x}^{1}: T_{x} M \rightarrow V$.



1-форма $\theta: T M_{1} \rightarrow V$, значение к-рой на векторе $X \in T_{u_{x}^{1}}^{M_{1}}$ равно $\theta_{u_{x}^{1}}(X)=u_{x}^{1}\left(\pi_{1}\right)_{*} X$, наз. ф о р м о й с м е щ е н и я. В локальных координатах $\left(x^{i}, u_{i}^{a}\right)$ многообразия $M_{1}$ форма $\theta$ имеет вид $\theta=u_{i}^{a} d x^{i} \otimes e_{a}$, где $e_{a}$ - стандартный базис в $V$.



Ограничение $\theta_{P}$ формы смещения $\theta$ на $G$-C. $P \subset M_{1}$ наз. ф о р м о й с м е щ е ни я $G-\mathrm{C}$. Она обладает следующими свойствами: 1) строгая горизонтальность: $\left.\theta_{P}(X)=0 \Leftrightarrow \pi_{*} X=0 ; 2\right) G$-әквивариантность: $\theta_{P} \circ=g \circ \theta_{P}$ для любого $g \in G$.



С помощью формы $\theta_{P}$ можно охарактеризовать главные расслоения с базой $M$, изоморфные $G$-С. А именно, главное $G$-расслоение $\pi: P \rightarrow M$ изоморфно $G-$.. тогда и только тогда, когда имеются точное линейное представление $\alpha$ групшы $G$ в $n$-мерном векторном пространстве $V, n=\operatorname{dim} M$, и $V$-значная строго горизонтальная $G$-эквивариантная 1 -форма $\theta$ на $P$. Отказ от требования точности представления $\alpha$ приводит к понятию обобщенной $G$-C. (1-го порядка) на $M$ - это главное $G$-расслоение $P \rightarrow M$ вместе с линейным представлением $\alpha: G \rightarrow$ $\rightarrow G L(V), \operatorname{dim} V=\operatorname{dim} M$, и $V$-значной строго горизонтальной $G$-әквивариантной 1 -формой $\theta$ на $P$.



Примером обобщенной $G-C$. является канонич. расслоение $\pi: P \rightarrow G \rightarrow P$ над однородным пространством $G \backslash P$ группы Ли $P$. Здесь $\alpha-$ изотропии представление



\includegraphics[max width=\textwidth, center]{2023_07_09_c51af5af36345b3c45c8g-1}

мой группы Ли $P$.



Пусть $\pi: P \rightarrow M$ есть G-С. 1-го порядка. Расслоение $\pi^{\prime}: P^{\prime} \rightarrow P$ 1-струй локальных сечений расслоения $\pi$ можно рассматривать как $G^{\prime}$-структуру на $P$, где $G^{\prime}=\operatorname{Hom}(V, \quad g)$ - коммутативная группа, $\mathfrak{g}-$ алгебра Ли групшы $G$, линейно представленная в пространстве $V \bigoplus \mathfrak{g}$ формулой



$A(v, X)=(v, X+A(v)), A \in G^{\prime}, v \in V, X \in \mathfrak{g}$,



и действующая на многообразии $P^{\prime}$ по формуле



$H \longmapsto A H=\left\{l_{p} A(\theta(h))+H, A \in G^{\prime}, p=\pi^{\prime}(H), h \in H\right\}$, где $l_{p}$ - канонич. изоморфизм алгебры Ли $\mathfrak{g}$ группы $G$ на вертикальное подпространство $T_{p}^{V} P=T_{p}\left(\pi^{-1}(\pi(p))\right)$. Элемент $H \in P^{\prime}$ может рассматриваться как горизонтальное (т. е. дополнительное к вертикальному) подпространство в $T_{p} P$. Он определяет корепер $\theta_{H}^{\prime}: T_{p} P \approx \mathfrak{g}+V$, к-рый на вертикальном подпростран- стве задается отображением $l_{p}$, а на горизонтальном отображением $\theta_{H}=\theta \mid H$. Вектор-функция $C^{\prime}: P^{\prime} \rightarrow W=$ $=\operatorname{Hom}(V \wedge V, V)$, заданная формулой $H \mapsto C_{H}^{\prime}, C_{H}^{\prime}(u, v)=$ $=d \theta\left(\theta_{H^{-1}}^{-1} u, \theta{ }_{H}^{-1} v\right), \quad$ наз. ф у н к ц и е й ч е н и я $G$-C. $\pi$. Сечение $s: x \mapsto H_{P(x)}$ расслоения $\pi \circ \pi^{\prime}: P^{\prime} \rightarrow M$ задает связность в расслоении $\pi$, а ограничение функции $C^{\prime}$ на $S(M)$ есть функция, задающая координаты тензора кручения әтой связности относительно поля кореперов $p(x)$.



Отображение $C^{\prime}: P^{\prime} \rightarrow W \quad G^{\prime}$-әквивариантно относительно вышеуказанного действия группы $G^{\prime}$ в $P$ и действия $G^{\prime}$ в $W$, задаваемого формулой



\[

A: w \longmapsto A w=w+\delta A

\]



где $\delta: G^{\prime} \rightarrow W,(\delta A)(u, v)=A(u) v-A(v) u$. Индуцированное отображением $C^{\prime}$ отображение $C: P \rightarrow G^{\prime} \backslash W$ наз. с т р Ук т у р н о й ф ун кци е й $G$-структуры $\pi$. Обращение в нуль функции $C$ равносильно существованию в расслоении $\pi$ связности без кручения.



Выбор подпространства $D \subset W$, дополнительного к $\delta G^{\prime}$, определяет подрасслоение $P^{(1)}=C^{\prime-1}(D)$ расслоения кореперов $\pi^{\prime}: P^{\prime} \rightarrow P$ со структурной группой $G^{(1)}=G^{\prime} \cap \mathrm{Ker} \delta \cong \mathfrak{q} \otimes V^{*} \cap V \otimes S^{2} V^{*} \subset V \otimes V^{* 2}, \quad$ т. $G^{(1)}$-структуру $\pi^{(1)}=\left.\pi^{\prime}\right|_{p^{(1)}}: P^{(1)} \rightarrow P$ на $P$. Она наз. пе р вым п род о лж е ни е м $G$-С. $\pi$. По индукции определяется $i$-е п р о д о л же ни е $\pi^{(i)}: p^{(i)} \rightarrow$ $\rightarrow P^{(i-1)}$ как $G^{(i)}$-структура на $P^{(i-1)}$, где группа $G^{(i)}$ изоморфна векторной группе $\mathfrak{g} \otimes S^{i} V^{*} \cap V \otimes S^{i+1} V^{*} \subset$ $\subset V \otimes V^{*(i+1)}$. Структурная функция $C^{(i)} i$-го продолжения наз. ст т ук т у рн ой ф у у кци и й $i$-го п о р я д к а $G-\dot{C}$. $\pi$.



Центральной проблемой теории G-С. является локальная п ро б ле ма ә к в и в а л е т т о с т инахождение необходимых и достаточных условий, при к-рых две $G$-С. $\pi: P \rightarrow M$ и $\bar{\pi}: \bar{P} \rightarrow \bar{M}$ с одной и той же структурной группой $G$ локально әквивалентны, т. е. существует локальный диффеоморфизм $\varphi: M \supset$ $\supset U \rightarrow \bar{U} \subset \bar{M}$ многообразий $M$ и $\bar{M}$, индуцирующий изоморфизм G-С. над окрестностями $U$ и $\bar{U}$. Частным случаем әтой проблемы является п р о б ле м а и н т ег р и р у е м о с т и - нахождение необходимых и достаточных условий локальной эквивалентности данной $G$-С. и стандартной плоской $G$-С. Проблему локальной әквивалентности можно переформулировать как проблему нахождения полной системы локальных инвариантов $G-\mathrm{C}$.



Для $O(n)$-структуры, к-рая отождествляется с римавовой метрикой, проблема интегрируемости была решена Б. Риманом (B. Riemann): необходимые и достаточные условия интегрируемости состоят в обращении в нуль тензора кривизны метрики, а локальная проблема әквивалентности - Э. Кристофффелем (Е. Christoffel) и P. Липшицем (R. Lipschitz): полная система локальных инвариантов римановой метрики состоит из ее тензора кривизны и его последовательных ковариантных производных (см. [1]).



Подход к решению проблемы әквивалентности основан на понятиях продолжения и структурной функции. $\mathrm{C}$ каждой $G-\mathrm{C}$. $\pi: P \rightarrow M$ 1-го порядка со структурной группой $G \subset G L(n)$ связываются последовательность продолжений



\[

\ldots \longrightarrow P^{(i)} \longrightarrow P^{(i-1)} \longrightarrow \ldots \longrightarrow P \stackrel{\pi}{\longrightarrow} M

\]



и последовательность структурных функций $C^{(i)}$. Для $O(n)$-структуры структурная функция $C^{(0)}=C$ на $P^{(0)}=P$ равна 0, а существенные части остальных структурных функций $C^{(i)}, i>0$, отождествляются с тензором кривизны соответствующей метрики и его последовательными ковариантными производными. Для интегрируемости G-С. $\pi$ необходимо и достаточно посто-





\end{document}